\chapter{Conclusões Parciais}
\label{cap:05}

A elaboração do protótipo Bicharada visa fomentar o desenvolvimento de jogos de quebra-cabeças, por meio da documentação detalhada: desde os estudos de critérios para bons jogos de quebra-cabeça até o rascunho e desenvolvimento de níveis. Busca-se também, um melhor entendimento das práticas de desenvolvimento de jogos e como certas decisões impactam a interação de jogadores com seu produto final, de forma positiva ou negativa, além de inspirar a criação de novos jogos de \textit{puzzle} e materiais acadêmicos sobre o tema. 

Dessa forma, almeja-se despertar interesse e engajamento em um público mais amplo, que pode não estar familiarizado com essa área específica da informática. Promovendo, discussões e reflexões que impulsionam, por meio de diretrizes e recomendações de boas práticas, novos \textit{designers} de jogos de \textit{puzzle}.

Como etapa subsequente, o projeto focará na implementação do Bicharada utilizando o motor de jogos Godot. Esta fase permitirá a convergência entre os rascunhos teóricos e o desenvolvimento prático na plataforma. Estima-se que a execução resulte em um protótipo funcional e estável, com suporte para os sistemas operacionais Windows e Linux.